\section{Introduction}
\label{sec:intro}

% Target: 0.75 page (~600 words). Write last — let the experiments shape the hook.

\todo{Open with a concrete hook: a figure (Fig.~\ref{fig:teaser}) showing a real
S\&P\,500 trajectory next to the synthetic series produced by the released
FTS-Diffusion artifacts. The contrast (real: rich pattern structure; synthetic:
near-linear ramp) is the visual thesis of the paper.}

% --- Background / problem ---
% Financial time-series generation is a high-value application of generative
% modeling — limited data history makes data augmentation appealing, but
% financial series have specific stylized facts (heavy tails, volatility
% clustering) that off-the-shelf generators struggle with.

% --- The paper under audit ---
% Huang et al.~\cite{huang2024ftsdiffusion} (ICLR~2024) propose FTS-Diffusion, a
% three-module generator (pattern recognition + diffusion-based segment
% generation + Markov-chain pattern evolution) and report (i) Kolmogorov–Smirnov
% / Anderson–Darling test statistics on returns that substantially exceed
% baselines (KS$\approx$0.32, AD$\approx$0.12 on three assets), and (ii)
% downstream MAPE reductions of up to 17.9\% on a one-day price-prediction task
% augmented with synthetic data.

% --- What we did ---
% We attempted to reproduce these results from the released code and
% configuration. The released sampling pipeline produces a degenerate
% near-linear generator: synthetic trajectories converge to a single tiled
% segment that violates every stylized fact the paper claims to preserve. We
% trace this to three specific implementation choices in the released code
% (Sec.~\ref{sec:reproducibility}). Surprisingly, several of the paper's
% evaluation metrics still pass on this degenerate output on the upward-trending
% S\&P\,500 asset (Sec.~\ref{sec:blindness}).

% --- Contributions ---
% \begin{itemize}
%   \item A reproducibility audit of FTS-Diffusion, with file:line evidence
%     of three implementation choices in the released code that cause the
%     degeneracy (Sec.~\ref{sec:reproducibility}).
%   \item Controlled experiments showing the paper's distribution-fit and
%     downstream-MAPE metrics fail to flag the degeneracy on the trending
%     asset — including a literal linear-ramp baseline that passes several
%     of the metrics (Sec.~\ref{sec:blindness}).
%   \item \todo{If §6 stretch is in:} A small stress-test framework
%     (pattern-distribution divergence, transition entropy, returns
%     second-moment decay) that catches the degeneracy and could have flagged
%     it pre-publication (Sec.~\ref{sec:stress}).
% \end{itemize}

% --- What we do not claim ---
% We do \emph{not} claim that the underlying method is structurally broken.
% Confirming that requires retraining the released code at the authors'
% specified epoch budget across multiple seeds (we report partial results in
% Sec.~\ref{sec:reproducibility}; see also Sec.~\ref{sec:discussion} for
% limitations). The narrower defensible claim — that the released artifacts as
% published do not reproduce the paper's results, and that the paper's metrics
% do not detect this — is what we make here.

\begin{figure}[t]
  \centering
  \todo{Figure 1: real S\&P\,500 (left) vs.\ synthetic from released
  FTS-Diffusion (right). Source: \texttt{figures/synthetic\_trajectories/sp500\_k14.png}
  or regenerate.}
  \caption{Real and synthetic S\&P\,500 price trajectories over 25{,}200
  trading days. The synthetic series produced by the released FTS-Diffusion
  artifacts is a near-linear ramp; the pattern structure visible in the real
  series is absent. \todo{Replace with final figure.}}
  \label{fig:teaser}
\end{figure}
